% Generated draft: requires mathematical and source review.
\documentclass[11pt]{article}
\usepackage[margin=25mm]{geometry}
\usepackage{fontspec}
\setmainfont{Latin Modern Roman}
\usepackage{amsmath,amssymb,mathtools,unicode-math}
\setmathfont{Latin Modern Math}
\usepackage{enumitem,needspace,xurl,hyperref}
\setlength{\parindent}{0pt}
\setlength{\parskip}{5pt}
\setlength{\emergencystretch}{3em}
\newcommand{\N}{\mathbb{N}}
\newcommand{\Z}{\mathbb{Z}}
\newcommand{\Q}{\mathbb{Q}}
\newcommand{\R}{\mathbb{R}}
\newcommand{\C}{\mathbb{C}}
\newcommand{\F}{\mathbb{F}}
\newcommand{\A}{\mathbb{A}}
\newcommand{\PP}{\mathbb{P}}
\newcommand{\E}{\mathbb{E}}
\newcommand{\eps}{\varepsilon}
\newcommand{\dd}{\,\mathrm{d}}
\begin{document}
\section*{UnsolvedMath ID: 30002056}
\textbf{Source problem code:} catalogue entry
\subsection{Definitions and mathematical statement}
For a positive integer \(n\) and \(p=p(n)\in[0,1]\), let \(G(n,p)\) be the random graph on \(\{1,\ldots,n\}\) in which each unordered pair of distinct vertices is an edge independently with probability \(p\). Under the standard interpretation of the supplied notation, \(X(n,p)\) is its flag complex: a vertex subset is a simplex precisely when every pair of its distinct vertices is an edge. Write \(|X(n,p)|\) for its geometric realization. An event holds asymptotically almost surely if its probability tends to \(1\) as \(n\to\infty\).

The source does not explicitly specify whether \(d\) must be an integer independent of \(n\). Under the customary interpretation that \(d\geq6\) is a fixed integer, interpret
\[
 n^{-2/d}\ll p(n)\ll n^{-2/(d+1)}
\]
as the two conditions
\[
 p(n)n^{2/d}\longrightarrow\infty,
 \qquad
 p(n)n^{2/(d+1)}\longrightarrow0.
\]
Set \(k=\lfloor d/2\rfloor\). For an integer \(m\geq1\), \(\bigvee_{j=1}^{m}S^k\) denotes the space obtained by identifying one chosen point of each of \(m\) copies of the \(k\)-dimensional sphere. Adopt the convention that the wedge with \(m=0\) is a point.

\textbf{Question.} For every fixed integer \(d\geq6\) and every sequence \(p(n)\in[0,1]\) satisfying the two displayed limits, is it true that
\[
 \lim_{n\to\infty}
 \Pr\!\left(
   \exists m\in\mathbb Z_{\geq0}:
   |X(n,p(n))|\simeq\bigvee_{j=1}^{m}S^{\lfloor d/2\rfloor}
 \right)=1?
\]
Here \(\simeq\) denotes homotopy equivalence, and the number \(m\) of spheres may depend on the realized complex. The source does not specify that number or whether the empty wedge is intended to be allowed; the displayed formulation uses the convention stated above.

The supplied background discusses an extracted torsion-vanishing conjecture with an unspecified \(p\)-range, whereas the supplied mathematical statement asks for a wedge-of-spheres homotopy type in an explicit \(p\)-range. Their intended relationship is unresolved in this record; the background does not provide a precise alternative torsion conjecture.
\subsection{Short English statement}
For a fixed integer \(d\geq6\), form the clique complex of a random graph whose edge probability satisfies \(n^{-2/d}\ll p\ll n^{-2/(d+1)}\). The question asks whether, with probability tending to one, this complex is homotopy equivalent to a finite wedge of spheres, all of dimension \(\lfloor d/2\rfloor\). The fixed-integer interpretation of \(d\), the empty-wedge convention, and the relationship to the background torsion discussion require review.
\subsection{Sources}
\begin{itemize}
\item Supplied local UnsolvedMath registry record: permanent ID \(30002056\), source problem code \texttt{catalogue entry}, title ``Wedge-of-Spheres Homotopy Type of Random Flag Complexes,'' in the Oberwolfach Reports Open Problems collection. The record gives proposed year \(2012\) and supplies no named proposer.
\item Supplied report citation: \emph{Triangulations (2013)}, Oberwolfach Reports, pp.~1405--1486, DOI \texttt{10.4171/owr/2012/24}. Source URL and literature-source label ``original OWR report'': \url{https://doi.org/10.4171/owr/2012/24}. The supplied citation's title year \(2013\) differs from the proposed year and the year appearing in the DOI; these details are retained without reconciliation.
\item Supplied literature source, labeled ``integer-homology threshold'': \url{https://doi.org/10.1007/s00454-017-9863-1}.
\item Supplied literature source, labeled ``random-complex survey'': \url{https://arxiv.org/abs/1301.7165}.
\item Supplied literature assessment, checked on \texttt{2026-08-21}: status ``partially solved.'' The assessment says that integer-homology vanishing thresholds are known in important random-complex models, but that the extracted torsion-vanishing conjecture has no specified \(p\)-range and the corresponding torsion-free behavior remains unresolved. This is supplied dated curation derived from the OWR collection, not an independently verified literature assessment; its time-sensitive status and cited claims require independent verification.
\item The registry labels the statement \texttt{corrected\_verified} and reports that it was checked against the official Oberwolfach report, preserving the stated hypotheses, notation, and open target. This is a supplied verification claim, not verification performed for this draft.
\end{itemize}
\end{document}
